\documentclass[11pt]{article}
%\usepackage[sort&compress,comma,numbers,round]{natbib}
\usepackage{graphicx}
\usepackage{bm}
\usepackage[hidelinks=true]{hyperref} 
\usepackage{xcolor}
\usepackage{amsmath}
\usepackage{amsfonts}
\usepackage{float, bigints}
\usepackage{amssymb,graphicx,fullpage}
\usepackage{mathrsfs}
\usepackage{bm}
\usepackage[hidelinks=true]{hyperref} 
\usepackage{mathrsfs, amsmath}
\usepackage{amsbsy,setspace}
\usepackage{authblk}

\usepackage{tcolorbox}
\tcbuselibrary{theorems}

\newtcbtheorem[]{mytheo}{}%
{colback=gray!5,colframe=gray!35!black,fonttitle=\bfseries}{th}

\renewcommand\Affilfont{\itshape\small}
%\definecolor{greenish}{rgb}{0.13,0.58,0.16}
\definecolor{greenish}{RGB}{141,198,63}
\definecolor{reddish}{RGB}{239,65,54}
\definecolor{blueish}{RGB}{28,117,188}

\hypersetup{
colorlinks   = true, %Colours links instead of ugly boxes
urlcolor     = blueish, %Colour for external hyperlinks
linkcolor    = magenta, %Colour of internal links
citecolor   =  greenish%Colour of citations
}

%% Math defs
\newcommand{\real}{{\mathbb{R}}}
\newcommand{\complex}{{\mathbb{C}}}
\newcommand{\realpositive}{{\mathbb{R}}_{>0}}
\newcommand{\realnonnegative}{{\mathbb{R}}_{\ge 0}}
\newcommand{\integernonnegative}{\mathbb{N}_0}
\newcommand{\naturals}{\mathbb{N}}
\newcommand{\ones}[1]{\mathbf{1}_{#1}}
\newcommand{\Sc}{\mathcal{S}}
\newcommand{\meas}[1]{|#1|}
\newcommand{\Xsp}{\mathbb{X}}
\newcommand{\Wsp}{\mathbb{W}}
\newcommand{\Ysp}{\mathbb{Y}}
\newcommand{\Vsp}{\mathbb{V}}
\newcommand{\yvec}{\mathbf{y}}
\newcommand{\crefrangeconjunction}{--}


%% Math operators
\newcommand{\dvol}{\operatorname{dvol}}
\newcommand{\dd}{\operatorname{d}}
\newcommand{\vol}{\operatorname{vol}}
\newcommand{\diff}{\operatorname{d}}
\newcommand{\hess}{\operatorname{Hess}}
\newcommand{\supp}{\operatorname{supp}}
\newcommand{\sgn}{\operatorname{sgn}}
\newcommand{\id}{\operatorname{id}}
\newcommand{\RMSE}{\operatorname{RMSE}}
\newcommand{\longthmtitle}[1]{\mbox{}{\textit{(#1).}}}
\newcommand{\Lip}{\operatorname{Lip}}

\newcommand{\x}{\operatorname{x}}
\newcommand{\ctrl}{\operatorname{u}}

\def\d{\mathrm{d}}
\def\e{\mathrm{e}}
\def\hp{\hat{p}}
\def\E{\mathbb{E}}
\title{Schr\"odinger Bridge Problem}
\author{}
\date{}

\begin{document}

\maketitle
\begin{singlespace}
\section{Introduction}
This primer is an absolute beginner's introduction to the Schr\"odinger Bridge Proglem (SBP), which has attracted a lot interest given the rise of generative algorithms (such as the diffusion model) inspired by this problem. Before we delve into the problem itself, I will quickly summarize a few key ideas from stochastic optimal control that are imperative for us to make progress.

\subsection{Stochastic differential equations}

Consider a Lagrangian particle undergoing arbitrary drift, $v(X_t, t)$ and diffusion motion, described by delta-correlated noise/Wiener process, $W_t$. Its dynamics is described by a stochastic differential equation (SDE) of the form,
\begin{equation}
    \d X_t = v(X_t, t) \ \d t + \sigma (X_t, t) \d W_t.
\end{equation}
Consider several instances of the dynamics of this particle with prescribed initial conditions $X(0) = X_o$, then we can describe the ensemble of these instances via a probability distribution, $p(x, t)$ which gives the probability of finding a particle at a particular position $x$ at time $t$. The evolution of this probability density is given by the Fokker-Planck equation,
\begin{equation}
    \partial_t p(x, t) + \partial_x[v(x, t) p(x, t)] = \partial_x^2 (D(x, t) p(x, t)),
    \label{eq:FPP}
\end{equation}
where $D(x, t) = \sigma^2(X_t, t)/2$.

Let us consider a simple example to make this concrete. If drift velocity is a constant, $v(X_t, t) = v_o$ and the amplitude of noise is also a constant, $\sigma(X_t, t) = \sqrt{2D}$, then the SDE is known as the Ornstein-Uhlenbeck process: $\d X_t = v_o \ \d t + \sqrt{2D} \d W_t$. The corresponding FP equation is given by,
\begin{equation}
    \partial_t p(x, t) + v_o\partial_x[p(x, t)] = D\partial_x^2 p(x, t).
    \label{eq:OUFPP}
\end{equation}
Initialising a particle at a position, $X_{t=0} = X_o$ corresponds to initialising the PDE with initial condition, $p(x, 0) = \delta(x-x_o)$. There are several ways of solving this PDE, most convenient is by rewriting the PDE in transformed variable, $z = x - v_o t$ to get,
\[
\partial_t p(z, t) = D \partial_z^2 p(z, t).
\]
This of course is the heat-equation and can be solved using various techniques, here we show using Fourier transform. Defining forward and inverse transforms as,
\begin{align}
    \hp(k, t) =&\ \int_{-\infty}^{\infty} p(x, t) \e^{ikz}\,\d z,\\
    p(z, t) =&\ \frac{1}{2\pi} \int_{-\infty}^{\infty} \hp(k, t) \e^{-ikz}\,\d k,
\end{align}
the heat equation reduces to an amplitude equation,
\[
\frac{\d \hp(k, t)}{\d t} = -D k^2 \hp(k, t),
\]
for every $k$-th mode. This can be solved easily to get, $\hp(k, t) = \hp_o(k) \exp(-Dk^2t)$. The initial condition $p(z, 0) = \delta(z-x_o)$ gives $\hp_o(k) = \exp(ikx_o)$. The final solution is formally,

\[
p(x, t) =  \frac{1}{2\pi} \int_{-\infty}^{\infty} \e^{-Dk^2 t - ik(z-x_o)} \ \d k.
\]
A very useful integral formula is,
\[
\int_\infty^\infty \e^{-a z^2 \pm b z}\,\d z = \sqrt{\frac{\pi}{a}} \e^{(b^2/4a)}.
\]
Using this, we can get the classical result,
\[
p(z, t) = \sqrt{\frac{\pi}{4Dt}}\, \e^{-(z - x_o)^2/4 D t}.
\]
In the original frame of reference, the solution reads,
\[
p(x, t) = \sqrt{\frac{\pi}{4Dt}}\, \e^{-(x - x_o - v_o t)^2/4 D t}.
\]
It is clear that the mean of the distribution, $\E[x] = (x_o+v_o t)$ while the variance,


% Consider an active particle whose motion is described by
% an overdamped Langevin equation
% \begin{align*}
%     \eta \dot{\x}(t) =& -\nabla \Psi(\x(t)) + \ctrl(t) + \sqrt{2D} \xi(t)
% \end{align*}
% where $\eta$ and $D$ are the damping and 
% diffusion constants respectively,
% $\Psi$ is a potential,
% $\ctrl$ is a control input (a driving force)
% and $\xi(t)$ is a delta-correlated, stationary Gaussian process.
% We set $\eta =1$ and express the dynamics (in It$\hat{\rm o}$ form) as 
% $d\x(t) = -\nabla \Psi(\x(t)) dt + \ctrl(t) dt + \sqrt{2D} dB_t$
% where $B_t$ is a standard Brownian motion. \\

% \noindent We are interested in steering the particle
% for a finite duration of time $[0,T]$
% across the landscape described by $\Psi$
% towards regions of low $\Psi$
% (e.g., $\Psi$ is a height function
% over the $(x,y)$-place and the landscape $(x,y,\Psi(x,y))$
% consists of peaks and valleys, and we want to 
% steer the particle towards the lowest valley. 
% Note that although the particle naturally tends to
% go downhill due to the force $-\nabla \Psi$
% even in the absence of the driving input $\ctrl$,
% the landscape $\Psi$ may be non-convex 
% and the particle may be trapped in an undesirable 
% valley, i.e., region of attraction of a local minimum of $\Psi$).
% Furthermore, steering the particle incurs a cost
% (which we assume to be quadratic in $\ctrl$),
% and we thereby seek to minimize (on average) the following cost
% \begin{align*}
%     \frac{\gamma}{2} \int_0^T \ctrl^2(t) dt + \Psi(\x(T)).
% \end{align*}
% Note that there is a likely trade-off between
% expending steering effort (first term) and 
% the potential reward of reaching low values of $\Psi$
% (second term). \\

% \noindent In order to formulate the problem more precisely,
% let $J(t, \x, \ctrl_{t:T})$ be the average cost 
% from time $t$ when the particle is known to be at $\x$ 
% at time $t$, under the control input 
% $\ctrl_{t:T}$ delivered over the remaining time $[t,T]$
% \begin{align} \label{eq:expected_cost}
%     J(t, \x, \ctrl_{t:T}) = \mathbb{E} \left[ \left. \frac{\gamma}{2} \int_t^T \ctrl^2(t') dt' + \Psi(\x(T)) ~\right|~  \x(t) = \x \right]
% \end{align}
% where the expectation $\mathbb{E}$ is taken over all trajectories
% generated by the dynamics of the particle under the control $\ctrl$
% satisfying the condition $\x(t) = \x$. The goal is then to minimize
% $J(t, \x, \ctrl_{t:T})$ with an appropriate choice of the control $\ctrl_{t:T}$ 
% \begin{align*}
%    V(t, \x) = \min_{\ctrl_{t:T}} J(t, \x, \ctrl_{t:T}).
% \end{align*}
% The function $V$ above is typically called the value function.

% \section{Deterministic optimal control ($D=0$ case)}
% When $D=0$, the equation of motion of the particle is given by
% $\dot{\x}(t) = -\nabla \Psi(\x(t)) + \ctrl(t)$. The value function then satisfies
% \begin{align} \label{eq:value_fnc}
% \begin{aligned}
%     V(t,\x) &= \min_{\ctrl_{t:T}} \left \lbrace \frac{\gamma}{2} \int_t^T \ctrl^2(t') dt' + \Psi(\x(T)) \right \rbrace \\
%     &= \min_{\ctrl_{t:T}} \left \lbrace \frac{\gamma}{2} \ctrl^2(t) dt + \int_{t+dt}^T \ctrl^2(t') dt' + \Psi(\x(T)) \right \rbrace \\
%     &= \min_{\ctrl(t)} \min_{\ctrl_{t+dt:T}} \left \lbrace \frac{\gamma}{2} \ctrl^2(t) dt + \int_{t+dt}^T \ctrl^2(t') dt' + \Psi(\x(T)) \right \rbrace \\
%     &= \min_{\ctrl(t)} \left[ \frac{\gamma}{2} \ctrl^2(t) dt + \min_{\ctrl_{t+dt:T}} \left \lbrace \int_{t+dt}^T \ctrl^2(t') dt' + \Psi(\x(T)) \right \rbrace \right]
% \end{aligned}
% \end{align}
% where we first merely decompose the cost incurred over the time horizon $[t,T]$
% into the cost at time $t$ and the cost over $[t+dt, T]$,
% which allows us to express the original minimization problem as a 
% nested minimization. The term $\min_{\ctrl_{t+dt:T}} \left \lbrace \int_{t+dt}^T \ctrl^2(t') dt' + \Psi(\x(T)) \right \rbrace$ can be seen to have the
% same structure as $V$ but with the time argument $t+dt$ and a possibly
% different position argument. Since under the control input $\ctrl(t)$ at time $t$ 
% the particle moves to a new position $\x(t+dt) = \x - \nabla \Psi(\x) dt + \ctrl(t) dt$, we can write
% \begin{align*}  
%     V(t, \x) = \min_{v} \left[ \frac{\gamma}{2} v^2 dt + V \left(t+dt, \x -\nabla \Psi(\x) dt + v dt \right) \right]
% \end{align*}
% and we obtain (by Taylor expansion to first order in $dt$)
% \begin{align*}
%     0 &= \min_{v} \left[ \frac{\gamma}{2} v^2 dt + V \left(t+dt, \x -\nabla \Psi(\x) dt + v dt \right) - V(t,\x) \right] \\
%     &= \min_{v} \left[ \frac{\gamma}{2} v^2 dt + \frac{\partial V}{\partial t} dt - \nabla V \cdot \nabla \Psi dt + \nabla V \cdot v dt \right] \\
%     &= \frac{\partial V}{\partial t} - \nabla V \cdot \nabla \Psi + \min_{v} \left \lbrace \frac{\gamma}{2} v^2 + \nabla V \cdot v \right \rbrace
% \end{align*}
% The minimizing $v$ above is given by $v^* = - \frac{1}{\gamma} \nabla V$
% and we obtain
% \begin{align} \label{eq:HJB}
%     \frac{\partial V}{\partial t} - \nabla V \cdot \nabla \Psi - \frac{1}{2\gamma} \left| \nabla V \right|^2 = 0.
% \end{align}
% Note that for time $t=T$, we get the boundary condition $V(T, \x) = \Psi(\x)$ (note from \eqref{eq:value_fnc} that the steering cost evaluates to zero, since the limits of the integral are equal). The above PDE is called the \emph{Hamilton-Jacobi-Bellman} (HJB) equation, and is solved backwards in time. \\

% \noindent In principle, the above approach involves solving the HJB PDE 
% globally in space and time, in order to implement the optimal control 
% $\ctrl^*(t, \x) = - \frac{1}{\gamma} \nabla V(t, \x)$ locally at $(t, \x)$.
% An alternative is to solve the optimal control problem at the level of
% individual trajectories, an approach we now outline below. \\

% \noindent Recall that the minimization in \eqref{eq:value_fnc},
% is a constrained minimization problem subject to the dynamics
% of the particle (starting from $t=0$ with the particle at $\x$)
% \begin{align*}
%     \min_{\ctrl_{0:T}} ~ \frac{\gamma}{2} \int_0^T \ctrl^2(t) dt + \Psi(\x(T))
%     \quad \text{subject~to} \quad  \dot{\x}(t) = -\nabla \Psi(\x(t)) + \ctrl(t),
%     ~ \x(0) = \x.
% \end{align*}
% The Lagrangian for the above problem, with multiplier $\lambda$
% for the ODE constraint, is given by
% \begin{align*}
%     \mathcal{L} = \frac{\gamma}{2} \int_0^T \ctrl^2(t) dt
%                     - \int_0^T \lambda(t) \cdot \left[ \dot{\x}(t) + \nabla \Psi(\x(t)) - \ctrl(t) \right] dt  + \Psi(\x(T))
% \end{align*}
% This allows us to express the original constrained 
% minimization problem as the nested, unconstrained optimization
% problem below
% \begin{align*}
%     \min_{\ctrl_{0:T}} ~\min_{\x_{0:T}} ~\max_{\lambda_{0:T}} ~\mathcal{L} (\x_{0:T}, \ctrl_{0:T}, \lambda_{0:T})    
% \end{align*}
% The removal of the ODE constraint makes $\x_{0:T}$ 
% a free trajectory (not subject to the dynamics 
% $\dot{\x}(t) = -\nabla \Psi(\x(t)) + \ctrl(t)$)
% in the Lagrangian.
% In first maximizing with respect to $\lambda_{0:T}$,
% the Lagrangian $\mathcal{L}$ evaluates to $+\infty$
% if the coefficient of $\lambda$ in the integrand is non-zero 
% at any time, and is zero otherwise.
% By nesting this maximization within the minimization
% with respect to $\x_{0:T}$, we force the minimization
% to select the latter condition by setting the
% coefficient of $\lambda$ in the integrand to $0$,
% which is a mere restatement of the ODE constraint
% and the $\min-\max$ optimization retrieves the
% original minimization problem in effect.
% The solution to the optimal control problem 
% is necessarily a critical point of the Lagrangian $\mathcal{L}$,
% given by (calculus of variations)
% \begin{align} \label{eq:pontryagin}
% \begin{aligned}
%     \dot{\x}^*(t) &= -\nabla \Psi(\x^*(t)) + \ctrl^*(t), \qquad \x^*(0) = \x \\
%     \dot{\lambda}^*(t) &= \nabla^2 \Psi(\x^*(t)) \lambda^*(t), \qquad \quad ~\lambda^*(T) = \nabla \Psi(\x^*(T)) \\
%     \ctrl^*(t) &= - \frac{1}{\gamma} \lambda^*(t)
% \end{aligned}
% \end{align}
% The final equation for the optimal control is
% obtained from minimization with respect to $\ctrl$
% in the Lagrangian optimization.
% We can now construct a function $\mathcal{H}
%  = \frac{\gamma}{2} \ctrl^2 + \lambda \cdot \left( -\nabla \Psi + \ctrl \right)$
% to rewrite the optimal solution as 
% $\dot{\x}^*(t) = \nabla_\lambda \mathcal{H}(\x^*(t), \ctrl^*(t), \lambda^*(t))$,
% $\dot{\lambda}^*(t) = - \nabla \mathcal{H}(\x^*(t), \ctrl^*(t), \lambda^*(t))$
% and $\ctrl^*(t) = \arg \min_v \mathcal{H}(\x^*(t), v, \lambda^*(t))$.
% Note that the system $(\x^*(t), \lambda^*(t))$ has the structure of
% a Hamiltonian dynamics (with $\mathcal{H}$ as the Hamiltonian)
% and the optimal control minimizes the Hamiltonian at every point along the
% optimal trajectory. This is the well-known \emph{Pontryagin Minimum Principle}. \\

% \noindent To highlight the connection between the Hamilton-Jacobi-Bellman 
% and Pontryagin formulations of the optimal control problem,
% we differentiate the HJB equation \eqref{eq:HJB} with respect to $\x$,
% to get
% \begin{align*}
%     \frac{\partial}{\partial t} \nabla V  
%             + \nabla^2 V \left( - \nabla \Psi - \frac{1}{\gamma} \nabla V \right)
%             - \nabla^2 \Psi \nabla V = 0 .
% \end{align*}
% The optimal trajectory $\x^*(t)$ 
% satisfies $\dot{\x}^*(t) = - \nabla \Psi(\x^*(t)) + \ctrl^*(t)
% = - \nabla \Psi(\x^*(t)) - \frac{1}{\gamma} \nabla V(t, \x^*(t))$,
% and we see that the sum of the first two terms in the PDE above
% is merely the derivative of $\nabla V$ along the optimal trajectory, i.e.,
% \begin{align*}
%     \frac{D}{D t} \nabla V (t, \x^*(t)) - \nabla^2 \Psi(\x^*(t)) \nabla V(t, \x^*(t)) = 0 .
% \end{align*}
% By identifying $\lambda^*(t) = \nabla V(t, \x^*(t))$ in the equation above,
% we get $\dot{\lambda}^*(t) = \nabla^2 \Psi(\x^*(t)) \lambda^*(t)$
% and the boundary condition $\lambda^*(T) = \nabla V(T, \x^*(T)) = \nabla \Psi(\x^*(T))$ (again from the boundary condition of \eqref{eq:HJB}),
% thereby retrieving the Pontryagin formulation \eqref{eq:pontryagin}.
% This allows us to interpret the HJB equation as the Eulerian formulation
% of optimal control and Pontryagin as the corresponding Lagrangian formulation
% (in the terminology of fluid mechanics).

% \section{Stochastic Optimal Control}
% In the stochastic case, we only develop the HJB formalism here.
% Recall the average control cost \eqref{eq:expected_cost}
% for the stochastic control problem.
% Following the same procedure as in the deterministic case,
% we see that the value function satisfies
% \begin{align*}
%     V(t,x) = \min_v \left \lbrace \frac{\gamma}{2} v^2 dt + \mathbb{E}\left[ V(t+dt, \x -\nabla \Psi(\x(t)) dt + \ctrl(t) dt + \sqrt{2D} dB_t ) \right] \right \rbrace
% \end{align*}
% and we obtain (by Taylor expansion to first order in $dt$)
% \begin{align*}
%     0 &= \min_v \left \lbrace \frac{\gamma}{2} v^2 dt 
%             +  \frac{\partial V}{\partial t} dt 
%                 -  \nabla V \cdot \nabla \Psi dt + \nabla V \cdot v dt + D \Delta V dt \right \rbrace \\
%     &= \frac{\partial V}{\partial t} + D \Delta V  -  \nabla V \cdot \nabla \Psi
%         + \min_v \left \lbrace \frac{\gamma}{2} v^2  + \nabla V \cdot v \right \rbrace.
% \end{align*}
% With the minimizer $v^* = - \frac{1}{\gamma} \nabla V$ above, we get
% \begin{align} \label{eq:stochastic_HJB}
%     \frac{\partial V}{\partial t} + D \Delta V  -  \nabla V \cdot \nabla \Psi - \frac{1}{2\gamma} \left| \nabla V \right|^2 = 0
% \end{align}
% and the boundary condition $V(T,\x) = \Psi(\x)$.
% Note that we recover the deterministic HJB equation \eqref{eq:HJB} 
% when $D=0$. Furthermore, under the transform $V(t,\x) = - \frac{1}{\beta} \log \varphi (t, \x)$
% (for some $\beta > 0$)
% the HJB equation \eqref{eq:stochastic_HJB} yields
% \begin{align*}
%     \frac{\partial \varphi}{\partial t} + D \Delta \varphi -  \nabla \varphi \cdot \nabla \Psi 
%             - \left( D - \frac{1}{2\gamma \beta}  \right) \frac{1}{\varphi} \left| \nabla \varphi \right|^2 = 0
% \end{align*}
% We now impose an Einstein relation $D = \frac{1}{2\gamma \beta}$
% to reduce the above to a parabolic equation
% \begin{align*}
%     \frac{\partial \varphi}{\partial t} + D \Delta \varphi -  \nabla \varphi \cdot \nabla \Psi  = 0
% \end{align*}
% and the boundary condition $\varphi(t, \x) = e^{-\beta \Psi(\x)}$,
% the solution to which can be obtained as a path integral
% by applying the Feynman-Kac formula
% \begin{align*}
%     \varphi(t, \x) = \mathbb{E} \left[ \left. e^{-\beta \Psi(\x(T))} ~\right|~  \x(t) = \x \right]
% \end{align*}
% where the expectation above is taken with respect to paths generated by
% the uncontrolled dynamics $d\x(t') = -\nabla \Psi(\x(t')) dt' + \sqrt{2D} dB_{t'}$
% satisfying the condition $\x(t) = \x$. We then obtain the 
% value function as
% \begin{align*}
%     V(t,\x) = - \frac{1}{\beta} \log \left( \mathbb{E} \left[ \left. e^{-\beta \Psi(\x(T))} ~\right|~  \x(t) = \x \right] \right)
% \end{align*}
% and the optimal control $\ctrl^*(t, \x) = - \frac{1}{\gamma} \nabla V(t, \x)$.
% The value function of the stochastic optimal control problem
% takes the form of a free energy wherein the (path integral) 
% partition function corresponds to a Boltzmann distribution
% over paths generated by the uncontrolled dynamics
% with the weight factor $e^{-\beta \Psi(\x(T))}$
% obtained from the cost of an uncontrolled path $\Psi(\x(T))$
% (note from \eqref{eq:expected_cost} that the uncontrolled dynamics
% only incurs the terminal cost $\Psi(\x(T))$).
% The interesting fact here is that the 
% value function for the optimal control can be
% computed entirely from the uncontrolled dynamics
% (we note that we have shown this property only 
% when the control $\ctrl$ enters linearly in the dynamics
% and the cost is quadratic in $\ctrl$).

% \section{Optimal control as Bayesian inference}
% We have seen that the value function $V(t,\x)$ satisfies
% \begin{align} \label{eq:value_fnc_diff_forms}
% \begin{aligned}
%     V(t, \x) &= \min_{u_{t:T}}~ \mathbb{E}_{\mathbb{Q}} \left[ \left. \frac{\gamma}{2} \int_t^T \ctrl^2(t') dt' + \Psi(\x(T)) ~\right|~  \x(t) = \x \right] \\
%     &= - \frac{1}{\beta} \log \left( \mathbb{E}_{\mathbb{P}} \left[ \left. e^{-\beta \Psi(\x(T))} ~\right|~  \x(t) = \x \right] \right)
% \end{aligned}
% \end{align} 
% where $\mathbb{P}$ is the probability distribution on the
% space of paths generated by the uncontrolled dynamics,
% whereas $\mathbb{Q}$ is the the probability distribution on the
% space of paths generated by the dynamics under the control
% $u_{t:T}$. \\

% \noindent Let $\frac{d\mathbb{P}}{d\mathbb{Q}} \sim e^{-\beta \frac{\gamma}{2} \int_t^T \ctrl^2(t') dt'}$ be the Radon-Nikodym derivative (informally, the ratio of probability densities of the
% distributions $P$ and $Q$), upto a normalizing factor (which we will assume to be $1$ here).
% We then get
% \begin{align} \label{eq:value_fnc_jensen}
% \begin{aligned}
%     V(t, \x) &= - \frac{1}{\beta} \log \left( \mathbb{E}_{\mathbb{P}} \left[ \left. e^{-\beta \Psi(\x(T))} ~\right|~  \x(t) = \x \right] \right) \\
%     &= - \frac{1}{\beta} \log \left( \mathbb{E}_{\mathbb{Q}} \left[ \left. e^{-\beta \Psi(\x(T))} \frac{d\mathbb{P}}{d\mathbb{Q}} ~\right|~  \x(t) = \x \right] \right) \\
%     &= - \frac{1}{\beta} \log \left( \mathbb{E}_{\mathbb{Q}} \left[ \left. e^{-\beta \left(\frac{\gamma}{2} \int_t^T \ctrl^2(t') dt' + \Psi(\x(T)) \right)} ~\right|~  \x(t) = \x \right] \right) \\
%     &\leq - \frac{1}{\beta}  \mathbb{E}_{\mathbb{Q}} \left[ \log \left( \left. e^{-\beta \left(\frac{\gamma}{2} \int_t^T \ctrl^2(t') dt' + \Psi(\x(T)) \right)} \right) ~\right|~  \x(t) = \x \right]  \\
%     &= \mathbb{E}_{\mathbb{Q}} \left[  \left. \left(\frac{\gamma}{2} \int_t^T \ctrl^2(t') + \Psi(\x(T)) \right)  ~\right|~  \x(t) = \x \right]
% \end{aligned}
% \end{align}
% where the inequality above is obtained from an application of Jensen's inequality.
% From the third equality above, we get
% \begin{align*}
%     e^{-\beta V(t,x)} = \mathbb{E}_{\mathbb{Q}} \left[ \left. e^{-\beta \left(\frac{\gamma}{2} \int_t^T \ctrl^2(t') dt' + \Psi(\x(T)) \right)} ~\right|~  \x(t) = \x \right]
% \end{align*}
% for any distribution $\mathbb{Q}$ of paths generated by 
% the controlled dynamics (with control $\ctrl$). 
% In particular, for any control $\ctrl$ that steers the particle
% from $\x$ at time $t=0$ to an $x^*$ at time $T$ 
% that satisfies $\Psi(x^*) = 0$, we get
% \begin{align*}
%      e^{-\beta (V(0,x) - V(0,x^*))} = \mathbb{E}_{\mathbb{Q}} \left[ \left. e^{-\beta \frac{\gamma}{2} \int_t^T \ctrl^2(t') dt' } ~\right|~  \x(0) = \x \right]
% \end{align*}
% We note that $\Delta V = V(0,x) - V(0,x^*)$ is the difference in free energy between states $x$ and $x^*$ and $W = \frac{\gamma}{2} \int_t^T \ctrl^2(t') dt'$ is the work done to steer the particle from $x$
% to $x^*$, and we obtain the well-known \textit{Jarzynski equality}
% \begin{align} \label{eq:Jarzynski}
%     e^{-\beta \Delta V} = \mathbb{E} \left[ e^{-\beta W} \right].
% \end{align}
% We now return to~\eqref{eq:value_fnc_jensen} 
% where we note from \eqref{eq:value_fnc_diff_forms} 
% that when we minimize the expectation on the final equality in~\eqref{eq:value_fnc_jensen} with respect to $\ctrl_{t:T}$ (or equivalently with respect to $\mathbb{Q}$,
% since the control $\ctrl_{t:T}$ induces the probability distribution $\mathbb{Q}$,
% and vice versa) we retrieve $V(t,\x)$ and the inequality becomes an equality.
% Therefore, we get
% \begin{align*}
%     V(t, \x) &= \min_{\mathbb{Q}} - \frac{1}{\beta}  \mathbb{E}_{\mathbb{Q}} \left[ \log \left( \left. e^{-\beta \left(\frac{\gamma}{2} \int_t^T \ctrl^2(t') + \Psi(\x(T)) \right)} \right) ~\right|~  \x(t) = \x \right] \\
%     &= \min_{\mathbb{Q}} - \frac{1}{\beta}  \mathbb{E}_{\mathbb{Q}} \left[ \log \left( \left. e^{-\beta \Psi(\x(T))} \frac{d\mathbb{P}}{d\mathbb{Q}} \right) ~\right|~  \x(t) = \x \right] \\
%     &= \min_{\mathbb{Q}} ~\mathbb{E}_{\mathbb{Q}}[\Psi(\x(T)] - \frac{1}{\beta} \mathbb{E}_{\mathbb{Q}} \left[ \log \left( \frac{d\mathbb{P}}{d\mathbb{Q}} \right) \right] \\
%     &= \min_{\mathbb{Q}} ~\underbrace{\mathbb{E}_{\mathbb{Q}}[\Psi(\x(T)]}_{\rm Average~terminal~cost} + \frac{1}{\beta} \underbrace{H(\mathbb{Q}|\mathbb{P})}_{\rm Relative~entropy}
%  \end{align*}
% The optimal control (minimizing $\mathbb{Q}$ above)
% minimizes a weighted combination of the KL-divergence (or relative entropy)
% to the uncontrolled distribution $\mathbb{P}$
% and the expected terminal cost. We note that the control
% cost $\frac{\gamma}{2} \int_t^T \ctrl^2(t') dt'$ 
% is given by the KL-divergence of $\mathbb{Q}$ to $\mathbb{P}$,
% encoded as the deviation from the uncontrolled distribution.
% In Bayesian terms, the uncontrolled distribution
% $\mathbb{P}$ is the prior, the optimal controlled distribution $\mathbb{Q}^*$
% is obtained as the posterior according to Bayes rule (solving for $\mathbb{Q}^*$
% above)
% \begin{align*}
%     d\mathbb{Q}^* = \frac{e^{-\beta \Psi(\x(T))} d\mathbb{P}}{\mathbb{E}_{\mathbb{P}}\left[ e^{-\beta \Psi(\x(T))} \right]}
% \end{align*}

% \nocite{RS:94, RS:86, HK:05, ET-ET:12, MD-SG:11, ET:09, CJ:97}

\end{singlespace}

\bibliographystyle{unsrt}
\bibliography{references}

\end{document}

